\selectlanguage{english}

\sheettitle
  {Solutions}
  {Consult only after a first complete written attempt}

\section*{Exercise 0 — The probability-space $L^p$ interface}

For $1\le p<\infty$,
\[
L^p(\mathbb P)
=
\left\{
X:\mathbb E[|X|^p]<\infty
\right\}\big/\!\sim_{\mathbb P},
\qquad
X\sim_{\mathbb P}Y
\Longleftrightarrow
X=Y\quad\mathbb P\text{-almost surely}.
\]
Both integrability and the equivalence relation are defined by $\mathbb P$; hence the ambient measure is part of the object.

The finite-measure embedding gives, for $1\le p<q\le\infty$,
\[
\|X\|_p
\le
\mathbb P(\Omega)^{1/p-1/q}\|X\|_q
=
\|X\|_q.
\]
Thus
\[
L^\infty(\mathbb P)
\subseteq
L^2(\mathbb P)
\subseteq
L^1(\mathbb P).
\]
The equality $\mathbb P(\Omega)=1$ is used when the finite-measure constant is reduced to one.

The statement $X\in L^p(\mathbb P)$ means that the $p$-th absolute moment $\mathbb E[|X|^p]$ is finite. The two product estimates are
\[
\|YZ\|_1
\le
\|Y\|_1\|Z\|_\infty
\]
by endpoint Hölder, and
\[
\|UV\|_1
\le
\|U\|_2\|V\|_2
\]
by Cauchy--Schwarz. Accordingly, $L^1$ controls finite expectations and finite signed RN numerators, $L^\infty$ supplies indicators and bounded tests, and $L^2$ controls conditional second moments and cross-products.

\section*{Exercise 1 — The finite-measure embedding}

Suppose first that $q<\infty$. Set
\[
r=\frac qp,
\qquad
r'=\frac q{q-p}.
\]
These exponents are conjugate. Hölder's inequality yields
\[
\begin{aligned}
\|f\|_p^p
&=
\int_E |f|^p\mathbf 1_E\,d\mu\\
&\le
\left(\int_E|f|^{pr}\,d\mu\right)^{1/r}
\left(\int_E\mathbf 1_E^{r'}\,d\mu\right)^{1/r'}\\
&=
\|f\|_q^p\mu(E)^{1-p/q}.
\end{aligned}
\]
Taking the $p$-th root gives
\[
\|f\|_p
\le
\mu(E)^{1/p-1/q}\|f\|_q.
\]
Finiteness of $\mu(E)$ is used exactly to ensure that $\mathbf 1_E\in L^{r'}(\mu)$.

If $q=\infty$, then
\[
|f|^p
\le
\|f\|_\infty^p
\quad\mu\text{-almost everywhere},
\]
so
\[
\|f\|_p
\le
\mu(E)^{1/p}\|f\|_\infty.
\]
For a probability measure the constants are one. Hence
\[
L^\infty(\mathbb P)
\subseteq
L^2(\mathbb P)
\subseteq
L^1(\mathbb P),
\qquad
\|X\|_1\le\|X\|_2\le\|X\|_\infty.
\]

\section*{Exercise 2 — From indicators to bounded tests}

If $Y\in L^1(\mu)$ and $Z\in L^\infty(\mu)$, then
\[
|YZ|
\le
\|Z\|_\infty |Y|
\quad\mu\text{-almost everywhere}.
\]
Integration proves that $YZ\in L^1(\mu)$ and
\[
\|YZ\|_1\le\|Y\|_1\|Z\|_\infty.
\]

Let
\[
Z=\sum_{k=1}^n a_k\mathbf 1_{A_k},
\qquad
A_k\in\mathcal G,
\]
be bounded and simple. By linearity and the assumed event identities,
\[
\int_EZ(f-g)\,d\mu
=
\sum_{k=1}^n a_k\int_{A_k}(f-g)\,d\mu
=0.
\]

Now let $Z$ be bounded and $\mathcal G$-measurable. Choose bounded $\mathcal G$-measurable simple functions $Z_n$ such that
\[
Z_n\longrightarrow Z
\quad\text{pointwise},
\qquad
|Z_n|\le\|Z\|_\infty.
\]
Then
\[
|Z_n(f-g)|
\le
\|Z\|_\infty |f-g|,
\]
and the right-hand side is integrable. Dominated convergence therefore gives
\[
\int_EZ(f-g)\,d\mu=0.
\]

If $f$ and $g$ are $\mathcal G$-measurable, put $u=f-g$ and take
\[
A_+=\{u>0\}\in\mathcal G,
\qquad
A_-=\{u<0\}\in\mathcal G.
\]
The event identities imply
\[
\int_{A_+}u\,d\mu=0,
\qquad
\int_{A_-}(-u)\,d\mu=0.
\]
Both integrands are nonnegative, hence they vanish almost everywhere. Therefore $u=0$ $\mu$-almost everywhere.

\section*{Exercise 3 — Density measures and change of measure}

Let $f\in L^1(\mu)$ and write $f=f^+-f^-$. Since
\[
\int_Ef^+\,d\mu<\infty,
\qquad
\int_Ef^-\,d\mu<\infty,
\]
the positive measures $f^+\mu$ and $f^-\mu$ are finite. Therefore
\[
\nu_f=f^+\mu-f^-\mu
\]
is a finite signed measure. Moreover, $\mu(A)=0$ forces both integrals over $A$ to vanish, so $\nu_f\ll\mu$.

Suppose that $\nu_f=\nu_g$ and put $u=f-g$. Then
\[
\int_Au\,d\mu=0
\qquad(A\in\mathcal A).
\]
Taking $A=\{u>0\}$ gives $\int u^+\,d\mu=0$, and taking $A=\{u<0\}$ gives $\int u^-\,d\mu=0$. Hence $u=0$ $\mu$-almost everywhere.

Finally, let $\nu=h\mu$ with $h\ge0$. The identity
\[
\int_E\varphi\,d\nu
=
\int_E\varphi h\,d\mu
\]
holds for indicators by the definition of $\nu$, for nonnegative simple functions by linearity, and for arbitrary nonnegative measurable functions by monotone convergence. For signed $\varphi$, apply the result to $\varphi^+$ and $\varphi^-$ provided
\[
\int_E|\varphi|\,d\nu
=
\int_E|\varphi|h\,d\mu
<\infty.
\]
Thus
\[
\varphi\in L^1(\nu)
\quad\Longleftrightarrow\quad
\int_E|\varphi|h\,d\mu<\infty.
\]

\section*{Exercise 4 — July RN consumer map}

\begin{enumerate}
  \item For the density generated by $f\in L^1(\mu)$, the numerator is the finite signed measure
  \[
  \nu_f(A)=\int_Af\,d\mu
  \]
  on $(E,\mathcal A)$; the reference is $\mu$, and the density is unique $\mu$-almost everywhere.

  \item For $B\in\mathcal F$ with $\mathbb P(B)>0$, the numerator is
  \[
  \mathbb P_B(A)
  =
  \frac{\mathbb P(A\cap B)}{\mathbb P(B)}
  \]
  on $(\Omega,\mathcal F)$; the reference is $\mathbb P$, and
  \[
  \frac{d\mathbb P_B}{d\mathbb P}
  =
  \frac{\mathbf 1_B}{\mathbb P(B)}
  \quad\mathbb P\text{-almost surely}.
  \]

  \item For an $\mathbb R^d$-valued $X$ satisfying $\mu_X\ll\lambda_d$, the numerator is
  \[
  \mu_X=X_\#\mathbb P
  \]
  on $(\mathbb R^d,\mathcal B(\mathbb R^d))$. The reference measure is $\lambda_d$. The resulting density is unique $\lambda_d$-almost everywhere.

  \item For $Y\in L^1(\mathbb P)$ and a sub-$\sigma$-field $\mathcal G\subseteq\mathcal F$, the numerator is
  \[
  \nu_Y^{\mathcal G}(A)
  =
  \mathbb E[Y\mathbf 1_A],
  \qquad A\in\mathcal G,
  \]
  on $(\Omega,\mathcal G)$; the reference is $\mathbb P|_{\mathcal G}$, and the conditional-expectation representative is unique $\mathbb P$-almost surely.
\end{enumerate}

The first two rows are explicit forward density constructions. The law-density row requires an RN converse after a dominating state-space measure has been chosen. The conditional-expectation row also uses an RN converse on $(\Omega,\mathcal G)$.

The pushforward $\mu_X$ is already a probability measure and therefore a law. It is not a density until a reference measure and an absolute-continuity relation have been specified.

For the signed conditional-expectation numerator, take its Jordan decomposition
\[
\nu_Y^{\mathcal G}
=
(\nu_Y^{\mathcal G})^+
-
(\nu_Y^{\mathcal G})^-.
\]
Both parts are finite positive measures absolutely continuous with respect to $\mathbb P|_{\mathcal G}$. Apply the positive RN theorem to each part and subtract the resulting densities.

\section*{Closed-book reconstruction}

\[
\boxed{Y\in L^q(\mathbb P),\ q>1}
\Longrightarrow
\boxed{Y\in L^1(\mathbb P)}
\]
is the finite-measure $L^p$ embedding. Next,
\[
Y\in L^1(\mathbb P)
\Longrightarrow
Y\mathbb P\text{ is a finite signed measure}
\]
is a forward density construction. A further RN representation is not automatic: one must first identify a reference measure, prove absolute continuity, and verify the hypotheses of the positive theorem or its legitimate signed extension. The representative is then unique only almost everywhere relative to that reference measure.
